\begin{proposition}[E.36 --- Proved: the two rational Kummer coordinates are $1/3$ and $1/2$]
\label{E-kummer-values}
Use the actual unit tangent for $q=Z/X$ at
$O=[1:-1:0]$, and the bracket and projector
normalizations of E.27--E.29. Put
\[
 v_6=\sum_{i=0}^2(e_{A_i}+e_{B_i}),
 \qquad v_2=e_{C_1}+e_{C_2}.
\]
For the genuine cubic action class $[\eta]$,
the complementary projection to the elliptic
component is
\[
 (1-e_{\mathcal A})[\eta]
   =\left[\left(
          \frac13v_6+\frac12v_2\right)\kappa_3(3),w\right]
 \quad\text{in }H^1(G_{\mathbb Q,\{3,\infty\}},
                         A_{\mathbb Q_3}(-1)).
 \tag{E.36.1}
\]
The bracket in this formula means the map
$H^1(V_{\mathbb Q_3})\to H^1(A_{\mathbb Q_3}(-1))$
induced by $v\mapsto[v,w]$.
Thus $(\lambda,\lambda,\mu)$ in the three
Tate coordinates of E.27 is exactly
\[
                  (1/3,1/3,1/2).
 \tag{E.36.2}
\]
Together with E.35, this expresses the full
rational class as these evaluated Kummer
terms plus the inverse image, under the
specified isomorphism $\Phi(-1)$, of
$-\iota_{w*}\operatorname{edge}_H
(\kappa(1-x)\smile\kappa(1-y))$.
The latter three-adic class and the integral
torsion remain to be evaluated.
\end{proposition}
\begin{proof}
We work first over $K=\mathbb Q(\zeta_3)$,
where every element of
$\mathcal A=E[3]\rtimes\mu_6$ is defined.
For $a\in\mathcal A$, choose a geometric
base-point path identifying its image of
the tangent at $O$ with our fixed base point.
The automorphism of the arithmetic fundamental
group supplied by the curve automorphism
then has the form
\[
 n\longmapsto\sigma_a(n),\qquad
 s_O(g)\longmapsto k_a(g)s_O(g).
 \tag{E.36.3}
\]
Here $\sigma_a$ is a continuous automorphism
of the geometric pro-three group, $k_a$ is
an actual continuous section cocycle, and
its section class is the image tangent
$a(s_O)$. The arithmetic homomorphism law
gives the exact identity
\[
 \sigma_a\rho_g\sigma_a^{-1}
        =\operatorname{Int}(k_a(g))\rho_g.
 \tag{E.36.4}
\]
Changing the chosen path conjugates these
maps geometrically. After abelianization,
$\alpha_a=[k_a^{\mathrm{ab}}]\in H^1(G_K,V_{\mathbb Q_3})$
is independent of that change.

Let $\omega$ be a lift of $w$ in the rational
degree-two-to-three extension defining $\eta$.
Every $a$ fixes $w$ on homology, so write
$\sigma_a\omega=\omega+t_a$ with $t_a\in A_{\mathbb Q_3}$.
On this extension, inner conjugation by
$k_a(g)$ adds $[k_a(g)^{\rm ab},w]$ to a lift
of $w$. Apply (E.36.4) to $\omega$ and use
$\rho_g\omega=\chi(g)(\omega+\eta_g)$.
The result, after dividing by $\chi(g)$, is
\[
 (a-1)\eta_g
    =\chi(g)^{-1}g(t_a)-t_a
                    +[k_a(g)^{\rm ab},w].
 \tag{E.36.5}
\]
Thus the first two terms form precisely the
coboundary in the twisted module $A(-1)$.
Taking classes and averaging the actual
$54$ identities gives
\[
 (1-e_{\mathcal A})[\eta]
       =-[\overline\alpha,w],
 \qquad
 \overline\alpha=\frac1{54}\sum_{a\in\mathcal A}\alpha_a.
 \tag{E.36.6}
\]
This proves the sign and identifies the
required coordinate with actual path classes.

Under the compact projection
$V_{\mathbb Q_3}\to H$, the class $\alpha_a$
is the Kummer class of $aO-O\in E[3]$.
It vanishes after tensoring with $\mathbb Q_3$.
The fixed rational splitting $V=H\oplus W$
therefore identifies $\overline\alpha$ with
a class in $H^1(G_K,W)$. This is a statement
about the cohomology class; the initially
chosen path cocycles need not have zero
compact coordinates as actual cochains.
The following unit pairings determine this
boundary class exactly.

For each cusp $c$ define its normalized
tangent-line equation by
\[
 \begin{aligned}
 L_{A_i}&=\zeta^{2i}Y+Z,&
 L_{B_i}&=\zeta^{2i}X+Z,&
 L_{C_i}&=\zeta^{2i}X+Y.
 \end{aligned}
\]
In particular $L_O=X+Y$. For $c\ne O$ the
unit $f_c=L_c/L_O$ on $U_K$ has divisor
$3(c-O)$. Its geometric Kummer linear form
on the boundary quotient is
\[
                   \ell_{f_c}=3(\delta_c-\delta_O).
 \tag{E.36.7}
\]
These eight linear forms are a basis of
$\operatorname{Hom}(W,\mathbb Q_3(1))$.
As in E.35 they vanish on the compact
summand. For a function with a zero or
pole at $O$, write $\operatorname{lt}_O$
for its first nonzero coefficient in the
specified parameter $q$.
Functoriality of the Kummer map applied to
the two sections in (E.36.3) gives
\[
 \ell_{f_c*}(\alpha_a)=
     \kappa_3\bigl(\operatorname{lt}_O(a^*f_c)\bigr)
       -\kappa_3\bigl(\operatorname{lt}_O(f_c)\bigr).
 \tag{E.36.8}
\]
Indeed the first coefficient in the
pullback is the leading value of $f_c$
on the transported tangent $a(s_O)$;
the second is its leading value on $s_O$.
Their difference is exactly the geometric
abelianized section cocycle under the
map induced by the unit $f_c$.

We evaluate the mean of (E.36.8).
The monomial automorphisms permute the nine
tangent lines up to scalars in $\mu_3$.
Every line occurs six times in each
product $\prod_{a\in\mathcal A}a^*L_c$.
Consequently
\[
       \prod_{a\in\mathcal A}a^*f_c\in\mu_3.
 \tag{E.36.9}
\]
The orders at $O$ in this product sum to
zero, so its leading coefficient is also
the product of the individual leading
coefficients. A root of unity has zero
rational three-adic Kummer class. Averaging
(E.36.8) therefore gives the exact equality
\[
 \ell_{f_c*}(\overline\alpha)
      =-\kappa_3\bigl(\operatorname{lt}_O(f_c)\bigr).
 \tag{E.36.10}
\]
The program \texttt{kummer\_coordinates.py}
checks every pullback of every tangent
line; for these normalizations all nine
ratio products in fact equal one. Only
their displayed root-of-unity property
is needed for (E.36.10).

On the branch through $O$,
$Y/X=-(1+q^3)^{1/3}$, so
\[
             L_O/X=-q^3/3+O(q^6).
\]
For $c\ne O$ this gives
\[
 \operatorname{lt}_O(f_c)=-3L_c(1,-1,0).
 \tag{E.36.11}
\]
For $A_i$ and $B_i$ these numbers are
$3\zeta^{2i}$ and $-3\zeta^{2i}$,
respectively. Their rational Kummer class
is $\kappa_3(3)$. For $C_1,C_2$ they are
$3(1-\zeta^{2i})$. The identities
\[
 (1-\zeta)^2=-3\zeta,\qquad
 1-\zeta^2=-\zeta^2(1-\zeta)
\]
show that their rational Kummer classes
are $\tfrac32\kappa_3(3)$.

To solve the resulting eight equations,
represent a boundary vector by its nine
cusp coefficients modulo the common
diagonal, and choose the representative
whose $O$ coefficient is zero. In these
coordinates (E.36.7) is three times the
$c$ coefficient. Equations
(E.36.10)--(E.36.11) thus give
\[
 \overline\alpha=
      -\left(\frac13v_6+\frac12v_2\right)\kappa_3(3)
             \quad\text{in }H^1(G_K,W).
 \tag{E.36.12}
\]
No information in another boundary direction
is discarded: the eight unit forms are a
basis and all their values have been computed.
The same program checks their eight rational
linear equations and records the resulting
vector in this exact quotient presentation.

Both $v_6$ and $v_2$ are stable under
$G_{\mathbb Q}$, with their Tate action,
and $\kappa_3(3)$ is defined over
$\mathbb Q$. The restriction map in
degree-one rational cohomology for the
quadratic extension $K/\mathbb Q$ is
injective by finite-group averaging. The inflation
map from the unramified arithmetic quotient
to $G_{\mathbb Q}$ is also injective in
degree one. Thus (E.36.6) and (E.36.12)
prove (E.36.1) over the original specified
arithmetic group. In the basis of E.27,
$v_6=b_0+b_1+b_2$ and $v_2=b_3$, giving
(E.36.2). The denominator three belongs
to the rational projector calculation;
it supplies no assertion about an
integral splitting or integral torsion.
E.28 eliminates the other projected
multidegree, and E.35 identifies the
remaining elliptic component, proving
the final assertion. The manuscript
anchor is \S51.16, folios 733v--736v,
scans 172--178, and LM7-E.23(i).
\end{proof}
